\section{Related Work}\label{sec:related-work}

\subsection{The Political Impacts of AI and Social Media}\label{sec:radicalization-polarization}
The transformative impacts of technology on political systems have long been studied by scholars in social computing, political science, and related fields. Today, social media platforms are recognized to play an increasing role in phenomena like \textit{polarization}, which is generally understood to be ``the process through which people extremize in their views through communication with like-minded others about their grievances,'' and \textit{radicalization}, ``a form of polarization that involves internalizing specific norms for [...] violent collective action'' \cite{Smithetal2025PrimerOnPoliticization}. Within literature on polarization, scholars further distinguish between \textit{ideological} or \textit{issue} polarization, which refers to the ``divergence in positions based on support for contrasting policies or issues,'' and \textit{affective} polarization, which refers to a state in which ``members of different social groups (or political parties) have increasingly negative feelings or attitudes towards outgroup/s'' \cite{Smithetal2024PolarizationCollectiveEngagement}. Researchers raise concerns that social media ranking algorithms may escalate political polarization by optimizing for engagement, which in turn rewards biased, sensationalizing, and/or misleading online content to the detriment of democratic society \cite{Strayetal2023AlgorithmicManagement}. Moreover, social media has also been identified as a novel vector for extremist groups to spread violent ideologies and form ties with potential recruits \cite{LeapAndYoung2021RadicalizationAndDeradicalization}.

The relationship between social technology and political systems has only been further complicated in recent years by the rapid development and uptake of AI technologies. While researchers have demonstrated that AI technologies have the potential to influence political decision-making \cite{Fisheretal2025BiasedLLMs}, their impacts on online polarization and radicalization remain unclear. Scholars like \citet{Hendricksetal2026LLMsExtremistAttitudes} have suggested that LLMs can fuel extremist attitudes by using universal moral framings in conversations with users. In contrast, \citet{SeviAndMekik2025AIWB} find in a different experiment that conversations with LLMs only weaken political attitudes, but do not strengthen them. Researchers have also begun proposing more \textit{active} AI-based interventions to reduce both ideological and affective polarization. For instance, \citet{Costelloetal2024ReducingConspiracyBeliefs} show that conversations with LLMs can be persuasive for reducing beliefs in conspiracy theories, \citet{Argyleetal2023LeveragingAI} use LLMs to promote feelings of being understood in political conversations by suggesting more respectful rephrasings of political messages, and \citet{Tessleretal2024HabermasMachine} show that LLMs can generate political statements that achieve consensus on topics of disagreement.

To date, however, many proposals for AI interventions targeted at promoting healthier online discourse have yet to deployed at scale on real-world platforms, having instead only been evaluated in controlled lab settings. To better understand the gap between lab and field settings, recent work from \citet{JungherrAndRauchfleisch2026HiddenCosts} elicits and compares reactions to human and AI-mediated forms of hypothetical political campaign outreach. They find that participants tend to evaluate AI-mediated outreach more negatively than human outreach, including being more likely to perceive AI-mediated outreach as a threat to freedom. These findings highlight the limited generalizability of prior work conducted in lab settings and raise the need for more democratic governance over applications of AI in online spaces.

An exception is Community Notes, a system developed by Twitter (now X) and launched publicly on Twitter in January 2021 \cite{Wojciketal2022Birdwatch}. Community Notes is a crowd-sourced fact-checking system that allows approved contributors to add context or corrections to potentially misleading tweets. Contributors can write notes on tweets they believe need clarification, but notes are initially not visible to the public. A note only becomes publicly visible if it receives ``helpful'' ratings from other contributors across a diversity of viewpoints, as determined by a machine learning algorithm. As a result, rather than relying on majority opinion, a tweet being ``noted'' requires consensus across ideological lines. The strictness of the consensus requirement, however, means that scaling Community Notes remains a persistent challenge: only 12.5\% of posts for which a note has been proposed actually have a note which has been rated helpful enough to be displayed \cite{Deetal2024Supernotes}, while social media bots like X's Grok provide appealing and immediate forms of fact-checking to users at scale \cite{ZhouAndHou2025FromNotesToBots}. More recently, in a three-month field evaluation study, \citet{LiAndBakker2026AIFactChecking} find that LLMs can successfully write community notes that achieve more positive ratings than human raters, indicating promising opportunities to enhance online discourse with AI.

\subsection{Deliberative vs. Agonistic Democracy}\label{sec:agonistic-democracy}
Prior proposals for using AI to promote healthier online discourse have largely operated under a framework of \textit{deliberative democracy}, either implicitly or explicitly. Typically understood to be exemplified by thinkers like John Rawls or J\"urgen Habermas, deliberative democracy views rational deliberation as playing a key role in the health of democratic societies \cite{Rawls1971TheoryOfJustice, Habermas1985TheoryOfCommunicativeAction}. In this model of democracy, citizens resolve political disagreements through respectful, free, and equal deliberation with each other. By yielding an acceptable consensus to all deliberative participants, this process is meant to confer legitimacy on its outcomes and secure more epistemically reliable decisions \cite{Landemore2017EpistemicTurn}. Evaluating a broad set of proposals for using LLMs to enhance democracy, for instance, \citet{LazarAndManuali2026LLMsDemocracy} find that such proposals generally fall into one of four use cases: (1) \textit{summarizing} public inputs to governmental decision-making processes; (2) generating widely acceptable statements that \textit{aggregate} diverse political opinions; (3) \textit{representing} individual preferences over unseen policy options; and (4) \textit{facilitating} democratic deliberation. Lazar and Manuali continue, however, that existing proposals hew ``as close to the Habermasian `ideal speech situation' as is practically feasible'' by placing participants in a ``privileged, jury-like position where they engage in unfettered, unmanipulated, dispassionate debate over a set of problems that do not, for the most part, directly concern them.'' They therefore echo the concerns raised above that these proposals may ultimately prove ineffective in settings which real democratic institutions must operate in, saturated with inequalities, deep ideological disagreements, and external social influences.

This objection against deliberative democracy is also taken up by \citet{Mouffe2000DemocraticParadox}, who defends an alternative theory of democracy she calls \textit{agonistic pluralism}. Like \citet{LazarAndManuali2026LLMsDemocracy}, Mouffe argues that rational deliberation cannot be a viable path to resolve political disagreement because the very act of ``postulating the availability of a public sphere where power would have been eliminated and where a rational consensus could be realized'' already eliminates the most essential elements of disagreement \cite{Mouffe2000DemocraticParadox}. Rather, for Mouffe the danger for democratic societies is that the framing of a political decision as a ``rational consensus'' becomes a way of precluding contestation of that decision while naturalizing existing hegemonic structures of power. The impossibility of eliminating power from the public sphere means that democratic societies must provide avenues to challenge power and reconstitute it in ways that are more compatible with democracy. To this end, Mouffe distinguishes between \textit{antagonistic} political relationships in which one's opponent is seen as an ``enemy to be destroyed,'' and \textit{agonistic} political relationships in which one's opponent is seen as an ``adversary [...] whose ideas we combat but whose right to defend those ideas we do not put into question'' \cite{Mouffe2000DemocraticParadox}. She argues that the challenge for democracy is to find ways of transforming antagonistic political relationships into agonistic ones, thereby creating legitimate ways for citizens to challenge political decisions without resulting in excessive conflict.

An increasing number of scholars have begun applying agonistic pluralism to issues in technological development and design. Among the earliest scholars to do so is \citet{DiSalvo2012AdversarialDesign}, who adopts Mouffe's notion of the ``adversary'' to explicate three tactics of ``adversarial design'': (1) \textit{revealing hegemony}, or using design to visualize relations of power in contested political domains; (2) \textit{reconfiguring the remainder}, or using design to challenge social assumptions underlying human-computer interactions; and (3) \textit{articulating collectives}, or structuring connections between computational artifacts in ways that allow people to participate in political action together. Building on DiSalvo's work, \citet{Crawford2016CanAnAlgorithmBeAgonistic} makes the case for understanding algorithms as agonistic sites which encode political logics and can be contested by publics.
% Algorithms like Wikipedia's ``view history'' mode, which makes backstage debates over content edits visible, exemplify an agonistic approach to algorithmic design that ``assume[s] perpetual conflict, recognize[s] contestation, and [...] promote[s] 'multiple constituencies honoring different moral sources.' '' \cite{Crawford2016CanAnAlgorithmBeAgonistic}.
Indeed, we contend that the emergence of communities like the subreddit r/GetNoted, which highlights cases of posters being ``roasted'' on Community Notes, reveals an inescapably agonistic remainder present even in explicitly consensus-based algorithms. More recent work from \citet{Angelietal2021UnsettlingPlay} and \citet{Shawetal2025AgonisticImageGeneration} find applying agonistic principles to the design of games and image generation interfaces, respectively, to be effective means of provoking users to critically reflect on preexisting assumptions. Also relevant is work from \citet{Sax2022AlgorithmicNewsDiversity}, who proposes that agonistic pluralism can enrich research on recommender systems by informing a more expansive conception of content diversity and bringing attention to the processes of negotiation that come with implementing recommender systems.

Here, we build on this body of work to question what it might mean to take an agonistic approach to de-antagonizing online political conversations. For this reason, our work differs from prior work which sees the reduction of ideological polarization \cite{Tessleretal2024HabermasMachine} or affective polarization \cite{Strayetal2023AlgorithmicManagement, Strayetal2026ProsocialRrankingChallenge} as a regulative ideal. For example, although we partially agree with \citet{Strayetal2023AlgorithmicManagement} that ideological polarization can play a key role in fostering ``constructive'' agonistic conflict, we disagree with their characterization of affective polarization as necessarily causing ``destructive'' antagonistic conflict. Agonistic democrats like \citet{Mouffe2000DemocraticParadox} have long emphasized ``the crucial role played by passions and affects in securing allegiance to democratic values.'' Mouffe observes that:
\begin{quote}
    ``When political frontiers become blurred, the dynamics of politics is obstructed and the constitution of distinctive political identities is hindered. Disaffection towards political parties sets in and it discourages participation in the political process. Alas, as we have begun to witness in many countries, the result is not a more mature, reconciled society without sharp divisions but the growth of other types of collective identities around religious, nationalist, or ethnic forms of identification.''
\end{quote}
Further evidence suggests that polarization may play a key role in motivating progressive forms of collective action \cite{Smithetal2024PolarizationCollectiveEngagement}, and conversely that attempts to restrict space for contestation can fuel violent radicalization when groups perceive that conventional political strategies are unable to succeed \cite{Smithetal2025PrimerOnPoliticization}. For this reason, \citet{Westphal2025AffectivePolarization} adapts Mouffe's definitions of antagonistic and agonistic relationships to distinguish between between antagonistic and agonistic affective polarization, where antagonistic affective polarization encompasses behaviors aimed at the elimination of outgroups such as calling for political violence, manipulating political procedures to the exclusion of opponents, and using dehumanizing language to refer to political opponents. Westphal concludes that \textbf{affective polarization need not be inherently harmful for democracy so long as adequate channels exist to transform antagonistic affective polarization into agonistic forms}, a view that we adopt below in designing and evaluating our interventions.

% \subsection{Automatic Fact-checking and Correction}\label{sec:fact-check}

% - 2 line history of fact checking literature
% - limitations of LLMs in fact-checking
% - Agentic AI and fact checking (deep-research like systems)
% - Current state of the art
% - Compare study design against the @grok phenomenon literature
% - - Collect the few best papers on the @grok / "ask the AI to fact-check" phenomenon
% - - Compare our design, researcher stance, and study design against their methods + findings

\subsection{Satire and Tone in Human-AI Interaction}\label{sec:satire-ai}
The influence of deliberative democratic theory further extends to constructions of tone in human-AI interactions. \citet{Chronis2026NLPAsLanguageIdeology} argues that natural language processing (NLP) research functions as ``language ideology'' that privileges certain forms of discourse as ``civil'' while denigrating others as ``toxic.'' In particular, dominant conceptions of civility in NLP research once again reflect ``an ideological model of political organization [...] rooted in the Habermasian ideal of the bourgeois public sphere'' which ``identifies civility with factual, descriptive language, prioritizing `substance' over form, `ideas and opinions' over emotions, and `rational argument' over casual or expressive speech.'' \cite{Chronis2026NLPAsLanguageIdeology}. The effect is that NLP research reinforces stereotypically white and male dialects of ``Standard English'' as an ideal, marking other dialects typically associated with non-white or non-male groups as deficient. Prior proposals for using AI to promote healthier online discourse similarly tend to align with hegemonic conceptions of civility, with works such as \citet{Deetal2024Supernotes} instructing LLMs to write fact-checks ``in a neutral and unbiased manner'' without expressing ``speculations or opinions,'' and \citet{Argyleetal2023LeveragingAI} using LLMs to suggest more polite rephrasings of political messages.

One form of communication often excluded by dominant conceptions of civil discourse is satire.
% Satire, derived from the Latin satura meaning a mixed dish,'' blends elements of irony, sarcasm, and exaggeration to communicate its perspective \cite{Holbert2015PoliticalSatire}.
\citet{Holbert2015PoliticalSatire} defines political satire as ``a pre-generic form of political discourse containing multiple humor elements that are utilized to attack and judge the flawed nature of human political activities.'' These four humor elements generally consist of: aggression, play, laughter, and judgment \cite{Holbert2015PoliticalSatire}. Scholars of humor like \citet{NorrickAndSpitz2008HumorAsAResource} have further documented a unique relationship between humor and social interaction. On one hand, forms of humor such as teasing, sarcasm, and ridicule may contain an aggressive or antagonistic dimension that carries out a social control function. On the other hand, humor may also have the ability to suspend conflict by lightening a conversation in the right situations. These qualities lead \citet{BoukesAndHameleers2023FactsOrHumor} to hypothesize that ``the humor in satire potentially weakens partisan-driven responses: opposed partisans are not necessarily perceived as the \textit{enemy} [...] but as \textit{legitimate discussion partners}.''\footnote{Emphasis ours.} Moreover, there is a long tradition of comedians using satire to ``punch up'' against existing social structures and build empathy with minoritized groups \cite{Mirowskietal2024LMsComedy}. For these reasons, satire may be a particularly apt method for agonistic democrats interested in creating conditions for agonistic forms of contestation between adversaries while defusing antagonistic forms of conflict between enemies.

Existing work on the relationship between AI, democracy, and humor is limited. Most relevant to our work is a study from \citet{BoukesAndHameleers2023FactsOrHumor}, who compare the effectiveness of ``regular'' and satirical fact-checks, respectively sourced from mainstream and satirical news websites. Importantly, they find that satirical fact-checks can be as effective as regular fact-checks at dissuading belief in misinformation, and may even be more effective for those who initially agree with the misinformation. However, they find that exposure to satirical fact-checks may increase affective polarization in those who initially agree with the fact-check. Scholars such as \citet{Mirowskietal2024LMsComedy} additionally investigate the competence of LLMs at generating humor through a series of workshops with comedians. They find that LLMs may be helpful at generating first drafts of jokes but are generally incapable of producing high-quality comedy due to a lack of perspective, connection with the audience, and contextual knowledge. To our knowledge, our study is the first to propose using satirical AI to de-antagonize online political conversations. Following \citet{Mirowskietal2024LMsComedy}'s findings, we do not purport that our bot is capable of producing \textit{genuine} satire, but only that it produces responses in a satirical \textit{tone}.